% Fuente: Auxiliar 10, «Conceptual».
Justifique su respuesta. Si la afirmación es falsa, entregue un contraejemplo.
\begin{enumerate}
    \item Si un punto $x^*$ cumple las condiciones necesarias de primer orden, entonces $x^*$ puede no ser un óptimo local.
    \item En un problema de optimización no convexo, todo mínimo local es también un mínimo global.
    \item Un problema de optimización no lineal convexo y acotado\footnote{Un problema es acotado inferiormente si su función objetivo no puede tomar valores arbitrariamente pequeños dentro del conjunto factible.}, siempre admite una solución óptima.
    \item Si un punto cumple con las condiciones de KKT para un problema no lineal, entonces necesariamente es una solución óptima.
\end{enumerate}
